\documentclass{article}
% Synthetic fixture for the four LLM-only claim-forward dispositions
% (latex-paper-en). Cases A-D are positives (pre-emptive rebuttal shell,
% concession that dismisses a confound, over-strong claim with a vague
% concession, empty modesty); Cases E-H are boundaries that must stay
% unchanged. Written for evals/evals.json id 25. The dispositions have no
% CF-* code, and the script tests do not read this file.
\begin{document}

\section{Introduction}
We study how a sample-reception protocol affects turnaround in a regional clinical laboratory network.

% Case A: pre-emptive rebuttal shell; it answers an objection that neither the data nor the cited work raises.
The reception protocol reduces median turnaround time from 96 to 71 minutes (Table~\ref{tab:turnaround}).
This does not mean that the protocol solves every delay in the laboratory.

\section{Data}
We interviewed 18 laboratory managers; Table~\ref{tab:interviews} lists their answers.
Eleven of the 18 managers named unclear reagent labelling as the main cause of sample rejection.
Table~\ref{tab:error} reports a labelling error of 4.2\% for night-shift samples and 2.1\% for day-shift samples.

% Case F: boundary; a definitional negation that keeps two quantities apart.
In this paper, turnaround time is the interval from sample receipt to result release; it does not include transport time.

\section{Results}
% Case B: concession that dismisses a confound; the denominator changed and no analysis is given.
Although monthly sample volume rose from 1,600 to 2,000 and the share of samples released on time fell from 90.0\% to 86.0\%, this does not affect the conclusion that the protocol improved on-time release.

% Case C: over-strong claim with a vague concession, on qualitative evidence.
The interviews conclusively prove that unclear reagent labelling is the decisive cause of sample rejection.
Admittedly, one cannot simply generalise this to every laboratory, but overall labelling plays a key role to some extent.

% Case E: boundary; a real evidence boundary that states its reason.
The interviews record the causes that managers report.
Because all 18 managers work in one regional network, the interviews do not measure how often each cause leads to a rejection elsewhere.

\section{Discussion}
% Case G: boundary; a concession about cited prior work, with the reason for the difference.
Although earlier audits~\cite{audit2019} report lower rejection rates, they counted only samples rejected at the bench, not samples returned at reception.

% Case D: empty modesty; no object, condition, or magnitude.
Our protocol still has certain limitations, and there is room for improvement.

\section{Conclusion}
The reception protocol reduced median turnaround time from 96 to 71 minutes across the network.

% Case H: boundary; a closing transition that names the concrete limitation, followed by the outlook.
Generalisation to multi-site networks remains limited because all interviews come from one regional network.
A multi-site survey with the same interview guide is the next step.

\end{document}
